\section{数据来源与变量构建}

\subsection{数据来源}

本文以福州市主城区为研究对象，覆盖鼓楼区、台江区、仓山区、晋安区与马尾区，必要时扩展至长乐区与闽侯县核心城镇片区。研究数据涵盖人口与社会经济、设施POI与空间交通三类。

\subsubsection{人口与社会经济数据}

人口空间分布数据采用国家地球系统科学数据中心发布的ASPECT数据集，该数据基于第七次全国人口普查乡镇级单元构建，空间分辨率100m，包含总人口、分年龄段人口与人口密度等栅格图层，坐标系为WGS 1984。本文以福州市行政边界裁剪汇总，提取各区老年人口规模与老龄化率等指标。社会经济数据包括地区生产总值、居民人均可支配收入与地方财政支出，来源于《福州统计年鉴》及各区统计公报。

\subsubsection{设施POI数据}

养老服务、医养协同与代际共享三类设施的兴趣点POI通过高德地图开放平台API关键词检索获取，每条记录含名称、经纬度、地址与所属行政区划等字段。养老类涵盖养老院、敬老院、护理院、日间照料中心、养老服务站、长者食堂与农村幸福院；医养协同类涵盖综合医院、社区卫生服务中心、药店与康复机构；代际共享类涵盖社区食堂、公园、图书馆、体育场馆、社区活动中心、便民服务中心、公交站与地铁站。

\subsubsection{空间交通数据}

路网与出行时间数据通过高德地图路径规划接口获取，公共交通站点经关键词检索获取，行政边界来源于全国地理信息资源目录服务系统，用于支撑可达性与服务圈覆盖范围测算。

\subsection{数据处理与空间匹配}

\subsubsection{空间网格划分}

为实现精细化空间分析，本文将研究区域划分为500m $\times$ 500m规则正方形网格，作为基本分析单元。街道尺度通常在数平方千米至数十平方千米，难以刻画社区级养老设施15至30分钟步行覆盖半径内的供需关系；500m网格面积为0.25平方千米，与老年群体日常活动半径相匹配，能更准确地度量服务可达性并识别补位点位。经划分，福州市全域共生成有效网格\_\_\_\_个。

\subsubsection{数据空间匹配}

将各类数据统一匹配至网格单元：人口栅格按网格边界分区统计，形成含总人口、老年人口与老龄化率的网格人口属性表；设施POI经坐标系统一、去重与类别校验后，按经纬度归入所在网格并分类汇总数量；公交、地铁站点同样归入网格，并利用路网出行时间计算网格中心至最近设施的通行耗时。

处理后，每个网格单元均具备人口特征、三类设施密度与交通可达性等属性，为后续建模奠定基础。

\subsection{指标体系设计}

\subsubsection{老年需求指标}

\begin{table}[H]
\centering
\caption{老年需求指标体系}
\begin{tabular}{ll}
\toprule
\rowcolor{tablehead}
指标 & 含义 \\
\midrule
老年人口密度 & 基础养老需求 \\
高龄人口密度 & 护理型养老需求 \\
人口密度 & 社区服务总体压力 \\
住宅小区密度 & 居住人口集聚程度 \\
老旧小区密度 & 潜在老龄人口集聚 \\
到医院距离 & 健康服务压力 \\
\bottomrule
\end{tabular}
\end{table}

说明：这一组指标回答“哪里更需要养老服务”。

\subsubsection{养老服务供给指标}

\begin{table}[H]
\centering
\caption{养老服务供给指标体系}
\begin{tabular}{ll}
\toprule
\rowcolor{tablehead}
指标 & 含义 \\
\midrule
养老机构数量 & 机构养老供给 \\
长者食堂数量 & 助餐服务供给 \\
日间照料中心数量 & 社区照护服务 \\
养老服务站数量 & 居家养老支持 \\
护理机构数量 & 失能照护供给 \\
养老床位数 & 机构容量，若缺失可用机构类型权重代替 \\
\bottomrule
\end{tabular}
\end{table}

说明：这一组指标回答“养老服务在哪里”。

\subsubsection{医养协同指标}

\begin{table}[H]
\centering
\caption{医养协同指标体系}
\begin{tabular}{ll}
\toprule
\rowcolor{tablehead}
指标 & 含义 \\
\midrule
医院密度 & 高等级医疗支持 \\
社区卫生服务中心密度 & 基层医疗支持 \\
药店密度 & 日常健康服务 \\
康复机构密度 & 康复照护支持 \\
长护险定点机构密度 & 护理型服务能力 \\
\bottomrule
\end{tabular}
\end{table}

说明：这一组指标回答“养老服务是否具备医疗支撑”。

\subsubsection{代际共享服务指标}

\begin{table}[H]
\centering
\caption{代际共享服务指标体系}
\begin{tabular}{ll}
\toprule
\rowcolor{tablehead}
指标 & 含义 \\
\midrule
社区食堂 / 平价餐饮 & 助餐与日常消费 \\
公园 / 广场 & 休闲活动 \\
图书馆 / 文化馆 & 公共文化服务 \\
体育空间 & 健康活动 \\
社区活动中心 & 社交互动 \\
便民服务中心 & 生活便利 \\
公交 / 地铁站 & 交通可达 \\
\bottomrule
\end{tabular}
\end{table}

说明：这一组指标回答“哪些服务能为养老压力提供间接缓冲”。
